\cotefile{19}
\scanpage{19}
3) Une suite

\[
0\to A'\unsure{-|\to} A\to A''-0 \tag*{(1)}
\]
dans $\mathcal C$ est dite $S$-exacte si la suite
\[
0-S(A')-S(A)\to S(A'')-0 \tag*{(1 bis)}
\]
\deleted{est exacte et} \inserted{est ``exacte et} \uline{se scinde}''.

Un objet $M\in\mathcal C$ est dit $S$-injectif, si pour toute suite $S$-exacte (1), la suite
\[
0-\mathrm{Hom}(A'',M)-\mathrm{Hom}(A,M)\longrightarrow\mathrm{Hom}(A',M)-0 \tag*{(2)}
\]
est exacte.

\uline{Exemples} (i) $M=T(X)$, car alors la suite (2) \unsure{s'écrira|s'écrit}
\[
0\to\mathrm{Hom}(S(A''),X)\to\mathrm{Hom}(S(A),X)-\mathrm{Hom}(S(A'),X)-0. \tag*{(2 bis)}
\]

[D'ailleurs, si (1) est telle que pour \deleted{tout} objet $S$-injectif $M$, \unsure{ou|on} \deleted{\unsure{particulièrement}} \inserted{\unsure{même}} seulement pour \unsure{tt} objet de la forme $M=\unsure{T|F|S}(X)$, \deleted{\unsure{avec} $X\in\mathcal C'$, (et même \unsure{seult} $X=S(A')$),} la suite (2) est exacte, alors (1) est $S$-exacte, car \deleted{(1 bis)} (2 bis) sera exacte pour \deleted{\unsure{tout} $X$ \illegible{2} pour $X=S(A')$,} \inserted{tt $X$, donc (1 bis) exacte}]

\deleted{\unsure{en} fait, il suffit de \unsure{prendre|fixer} $X$ \unsure{pour} \illegible{3} \unsure{Homomorphisme identique} de $S(A')$, et si $S$ est un foncteur exact à droite, alors il suffit pour $A'\to A-A''-0$ soit exacte, et que la suite}

(ii) Supposons $S$ \uline{exact} \inserted{à droite} \deleted{et ``fidèle'' en ce sens que} $S(u)=0\Longrightarrow u=0$. Soit $A\in\mathcal C$, je dis que la suite

\[
\deleted{\unsure{\Sigma}}\,\Sigma(A):\quad 0-A\to TS(A)\longrightarrowZ(A)-0 \tag*{(3)}
\]

(où $Z(X)=\mathrm{Coker}(X-TS(X))$) est $\unsure{S|\Sigma}$-exacte. En effet, \deleted{elle} se transforme par $S$

% archive sheet number: 18
